\scanpage{289}
Soit $Y=\operatorname{Spec}(V)$, $V$ anneau de valuation discrète\\
$X$ schéma abélien sur $Y$\\
$y_0,y_1$ p\textsuperscript{t} fermé resp. gén. de $Y$\\
$X_0,X_1$ les deux fibres de $X$ \ldots\\
$F,G$ deux sous-groupes de $X$, \uline{finis} sur $Y$,

\[
\begin{aligned}
\text{a)}\quad &F_1\not\subset G_1\ \deleted{\unsure{\ne0}}\\
\text{(ii)}\quad &F_0\subset G_0\ne0\\
\text{\deleted{\unsure{(ii)|(iii)}}\inserted{(iii)}}\quad &pF_{\unsure{0|1}}=G_1\qquad(p=\operatorname{car} k(V))
\end{aligned}
\]

\[
\begin{aligned}
A=X/F &\qquad X\xrightarrow{\alpha}A\\
B=X/G &\qquad X\xrightarrow{\beta}B
\end{aligned}
\qquad \text{hom can.}
\]

Considérons $\beta:X\mathbin{-}X/G$ (canonique) alors $p\beta$ s'annule sur \deleted{$pX$, donc sur} $F$\inserted{ en vertu de (iii)}, donc se factorise \uline{par} $u:X/F\mathbin{-}X/G$

\diagram{\begin{tikzcd}[ampersand replacement=\&,row sep=small,column sep=small]
\& X\arrow[dl,"\alpha"']\arrow[dr,"p\beta"] \&\\
A\arrow[rr,"u"] \& \& B
\end{tikzcd}}

\deleted{Cependant} $u_1\notin p\operatorname{Hom}(A_1,B_1)$, en \deleted{\unsure{eff}} se plaçant dans $\operatorname{Hom}(X_1,X_1)\otimes\mathbb{Q}\simeq\operatorname{Hom}(A_1,B_1)\otimes\mathbb{Q}$, si on avait $u=pv$, on aurait $v=\delta$, or $\delta$ ne provient pas d'un élément de $\operatorname{Hom}(A_1,B_1)$, puisque $F_1\not\subset G_1$ \deleted{[sinon, par (i), on aurait $F_1=0$ d'où $F=0$, ce qui contredit (ii)]}. Cependant $u_0\in p\operatorname{Hom}(A_0,F_0)$, puisque $F_0\subset G_0$

Les conditions dites (i)(ii)(iii) peuvent être réalisées, en égales ou inégales car, avec pour $X$ une \uline{courbe elliptique}, en prenant $X_0$ sans pt d'ordre $p$, $X_1$ avec pt d'ordre $p$, et $F_1$ et $G_1$ deux ss-groupes de $X_1$, \deleted{annulés par $p$}\inserted{de \deleted{\unsure{ord|ordre}} rang $p$ sur $k(y_1)$ et} distincts. Comme il n'y a

% archive sheet number: 288

